\section*{Chapitre \ref{ch:s1:positioning-and-sentiment} --- Positionnement et sentiment}

\begin{solution}{exo:s1:positioning-and-sentiment:1}
À partir de la clôture du vendredi (après la publication de 15~h~30), ou de l'ouverture du lundi par prudence~; plus tard pendant les semaines fériées. Utiliser le signal le mardi donne au backtest trois jours
d'information que personne ne possédait, et lui attribue la partie la plus fraîche et la plus précieuse du signal.
\end{solution}

\begin{solution}{exo:s1:positioning-and-sentiment:2}
$\phi = 2^{-1/20}$~: $\phi^3 = 0.901$, $\phi^{10} = 0.707$, $\phi^{21} = 0.483$. Les CI mesurés diminuent un peu plus vite (à 86~\%, 66~\% et 38~\% de
celui obtenu immédiatement), car un rapport sert pendant cinq jours après sa publication.
\end{solution}

\begin{solution}{exo:s1:positioning-and-sentiment:3}
Chaque contrat à terme comporte une position longue et une position courte. Les catégories du rapport couvrent presque toute la position ouverte, et pour le maïs les deux plus fortes positions nettes
sont celles des couvreurs (producteurs et négociants, nettement vendeurs) et celles des spéculateurs (gestion pour compte de tiers)~: quand l'une s'alourdit, l'autre prend la contrepartie.
\end{solution}

\begin{solution}{exo:s1:positioning-and-sentiment:4}
Environ $1.96/\sqrt{113} = 0.184$. Le 0.098 du maïs en vaut à peu près la moitié ($t = 1.04$).
\end{solution}

\begin{solution}{exo:s1:positioning-and-sentiment:5}
Une partie de la position des spéculateurs relève de la poursuite de tendance, et les tendances insérées dans \texttt{firm.synthfut} ont une demi-vie d'un an~: cette partie du signal
prédit les rendements plus longtemps que la \omterm{def:s1:positioning-and-sentiment:hedging}{pression de couverture}, dont la demi-vie est de 20 jours.
\end{solution}

\begin{solution}{exo:s1:positioning-and-sentiment:6}
Dans le modèle, les spéculateurs sont rémunérés pour absorber la couverture, si bien que les fortes positions longues spéculatives s'accompagnent de rendements espérés élevés~: aller à contre-courant de ces positions fait perdre de l'argent. Une règle à contre-courant
suppose que ce soient les spéculateurs qui paient~: des positions qui écartent les prix de leur valeur puis se débouclent, un mécanisme d'encombrement que le modèle ne contient pas.
\end{solution}

\begin{solution}{exo:s1:positioning-and-sentiment:7}
Le CI passe de 0.032, connu immédiatement, à 0.019, connu le vendredi~: le retard coûte 42~\% du signal, contre 14~\% avec une demi-vie de vingt jours. Plus
la pression dure peu de temps, plus un retard de trois jours coûte cher.
\end{solution}

\begin{solution}{exo:s1:positioning-and-sentiment:8}
Des valeurs d'enquête datées de leur période de référence et non de leur publication donnent au backtest une information antérieure à sa diffusion publique~; des indices construits
par composantes principales sur l'ensemble de l'échantillon utilisent des données postérieures pour pondérer des valeurs antérieures~; et les proxies ont été choisis a posteriori. Il faut utiliser les dates de publication
et un indice reconstruit avec les données disponibles à chaque date.
\end{solution}

\begin{solution}{pb:s1:positioning-and-sentiment:1}
\begin{enumerate}
\item En général le vendredi à 15~h~30, heure de l'Est américain, pour le mardi précédent.
\item La position nette des couvreurs rapportée à la position ouverte~; un signal issu des positions déclarées, telles qu'elles étaient connues à la date.
\item Les rendements des contrats à terme sur devises et sur produits agricoles varient avec les avoirs nets des couvreurs~; la \omterm{def:s1:positioning-and-sentiment:hedging}{pression de couverture} propre et celle du groupe influent sur les rendements dans 20 marchés.
\item Comme des images miroir~: une corrélation de $-0.97$.
\item Le score z sur 52 semaines de la gestion pour compte de tiers à chaque fin de mois, comparé à la variation du prix du maïs sur le deuxième mois suivant, de sorte que le mois du signal ne
chevauche pas le rendement.
\item 0.098, $t = 1.04$.
\item $+4.6~\%$ sur trois mois après des positions longues record, $+0.2~\%$ après des positions courtes record, $-0.9~\%$ sinon.
\item Rien dans un sens comme dans l'autre.
\item Un AR(1) de demi-vie de 20 jours rapportant 0.3 volatilité par an et par écart type~; les spéculateurs le détiennent, auxquels s'ajoutent la poursuite de tendance et le bruit.
\item 0.044, 0.038, 0.029 et 0.017 pour des retards de 0, 3, 10 et 21 jours.
\item Trois jours conservent $\phi^3 = 0.90$ de la pression~; un rapport utilisé pendant cinq jours perd un peu plus~: 14~\%.
\item Dans le modèle, les spéculateurs sont rémunérés, si bien que leurs extrêmes s'accompagnent de rendements espérés élevés.
\item \textbf{Résultat nommé.} Le \omterm{def:s1:positioning-and-sentiment:hedging}{signal de positionnement} des couvreurs a un IC de rang de 0.044 avec les rendements des cinq jours suivants lorsqu'il est connu immédiatement, et de 0.038
lorsqu'il est publié trois jours plus tard (14~\% de moins), tombant à 0.017 à 21 jours~; celui des spéculateurs, 0.043 et 0.039.
\item Une synthèse de l'optimisme des investisseurs tirée d'enquêtes ou de proxies de marché~; un faible sentiment est suivi de rendements relativement élevés sur les actions difficiles à valoriser.
\item Comme variable de conditionnement lente, qui indique quelles actions sont les plus mal évaluées, et non comme signal de négociation hebdomadaire.
\item Les positions de ceux qui paient pour la convexité, quotidiennes et par prix d'exercice.
\item Les rapports seulement à partir de leur date de publication, décalages des jours fériés compris, ainsi que les coûts.
\item Le contrarien des positions extrêmes.
\item Les deux rémunèrent celui qui porte un risque que d'autres veulent céder~: le carry le long de la courbe, la \omterm{def:s1:positioning-and-sentiment:hedging}{pression de couverture} au travers des positions.
\item Les couvreurs qui paient pour se défaire du risque de prix.
\end{enumerate}
\end{solution}

\begin{solution}{iq:s1:positioning-and-sentiment:1}
La position ouverte et les positions longues et courtes des catégories de négociants déclarables (producteurs et négociants, teneurs de swaps, gestion pour compte de tiers, autres), chaque semaine
pour le mardi~; il ne permet pas de connaître les positions prises dans la semaine, celles sur d'autres marchés ou de gré à gré, ni pourquoi une catégorie négocie.
\end{solution}

\begin{solution}{iq:s1:positioning-and-sentiment:2}
Les couvreurs paient pour transférer le risque de prix~: les producteurs qui vendent des contrats à terme en dessous du prix spot espéré laissent une prime aux spéculateurs qui les achètent.
\end{solution}

\begin{solution}{iq:s1:positioning-and-sentiment:3}
Une position longue record ne constitue pas à elle seule un signal~: il faut vérifier si elle absorbe de la couverture (positions des producteurs, forme de la courbe, stocks) ou si elle relève d'un encombrement (crowding) sur une tendance~;
il faut vérifier les preuves disponibles pour ce marché~; il faut dimensionner petit toute position à contre-courant, car les extrêmes peuvent persister et, dans certaines données, être suivis de gains supplémentaires.
\end{solution}

\begin{solution}{iq:s1:positioning-and-sentiment:4}
Chaque rapport avec sa date de référence (le mardi) et son horodatage de publication, issu du calendrier de publication de la CFTC~; les révisions conservées comme de nouvelles versions~; les requêtes
toujours classées par heure de publication.
\end{solution}

\begin{solution}{iq:s1:positioning-and-sentiment:5}
Choisir des proxies motivés (enquêtes, émissions, flux, rendements des introductions en bourse), les orthogonaliser par rapport aux conditions macroéconomiques, les combiner avec des pondérations
estimées sur les seules données passées, et tester l'indice sur la coupe transversale des actions difficiles à valoriser, hors échantillon (out-of-sample).
\end{solution}

\begin{solution}{iq:s1:positioning-and-sentiment:6}
Avec $r_{t+1} = a s_t + e_{t+1}$, $r_{t+k} = a s_{t+k-1} + e_{t+k}$ et $s_{t+k-1} = \phi^{k-1}s_t + \text{termes indépendants}$, la corrélation de $s_t$
avec $r_{t+k}$ vaut $\rho\phi^{k-1}$. Un rapport utilisé les jours $t + \ell, \dots, t + \ell + 4$ a une corrélation moyenne de $\rho\phi^{\ell}(1 - \phi^5)/(5(1 - \phi))$.
\end{solution}
